\subsection*{Conclusión individual --- Heiner Ariel Zambrano Ronquillo}
\addcontentsline{toc}{subsection}{Conclusión individual --- Heiner Ariel Zambrano Ronquillo}

El diseño del núcleo común (\texttt{lamport.py}, \texttt{config.py} y \texttt{machine\_info.py}), del orquestador experimental (\texttt{run\_experiments.py}) y del verificador causal (\texttt{verify\_lamport.py}) permitió comprobar empíricamente que comparar protocolos distribuidos exige rigor metodológico y no solo pruebas aisladas. Sobre el equipo controlado (\texttt{DESKTOP-OSLA9DF}, Windows 11, Ryzen 5 7535HS), descartar la corrida de calentamiento y repetir cinco veces cada escenario reveló que Sockets TCP con conexiones persistentes y \texttt{TCP\_NODELAY} superó a gRPC en latencia media en los cuatro casos: 0,0884\,ms frente a 0,4655\,ms en 100\,B secuencial ($+$426,77\,\% en gRPC) y 0,1988\,ms frente a 0,5774\,ms en 10\,KB concurrente ($+$190,39\,\%). Sin embargo, gRPC mantuvo su rendimiento concurrente estable al pasar de 100\,B a 10\,KB (5\,153,6 frente a 5\,151,1\,msg/s) y redujo la brecha relativa a menos de la mitad bajo carga pesada simultánea.

Por otro lado, auditar los 40 eventos reales por mecanismo confirmó un 100\,\% de cumplimiento de $L(\mathrm{recv}) > L(\mathrm{send})$ ($20/20$, $L_{\mathrm{max}}=38$) y evidenció pares concurrentes reales con idéntico reloj lógico (\texttt{ex-01} en \texttt{nodo-A} y \texttt{ex-02} en \texttt{nodo-B} con $L=1$), demostrando en práctica la necesidad de relojes vectoriales cuando se requiere detectar causalidad exacta. Trasladado a nuestro PFC (\emph{Bolsa de Empleos UTEQ}), esta experiencia ratifica que la elección tecnológica debe guiarse por contratos tipados, mantenibilidad y requisitos CAP (\emph{CP} con gRPC en validaciones síncronas internas y colas \emph{AP} en notificaciones) más que por microsegundos ganados en \texttt{localhost}.
